\documentclass[12pt,a4paper]{article}

% --- PAQUETES ---
\usepackage[utf8]{inputenc}
\usepackage[T1]{fontenc}
\usepackage[spanish,es-nodecimaldot]{babel} % Español, punto decimal
\usepackage{amsmath, amssymb, amsthm, mathtools} % Matemáticas avanzadas
\usepackage{geometry}
\usepackage{xcolor}
\usepackage{hyperref}
\usepackage{physics}

% --- CONFIGURACIÓN DE MÁRGENES ---
\geometry{
 a4paper,
 left=25mm,
 right=25mm,
 top=25mm,
 bottom=25mm
}

% --- ENTORNOS TIPO TEOREMA ---
\newtheorem{teorema}{Teorema}
\newtheorem{proposicion}[teorema]{Proposición}
\newtheorem{lema}[teorema]{Lema}
\newtheorem{definicion}[teorema]{Definición}
\newtheorem{observacion}[teorema]{Observación}

% --- COMANDOS PERSONALIZADOS ---
\newcommand{\C}{\mathbb{C}}
\newcommand{\R}{\mathbb{R}}
\newcommand{\Poly}{\mathbb{P}[z]}
%\newcommand{\ip}[2]{\langle #1, #2 \rangle}
\newcommand{\D}{\mathcal{D}}

% --- TÍTULO ---
\title{Transformaciones de ceros}
\date{}
\author{}

\begin{document}

\maketitle

Sea $\Poly$ el espacio de polinomios con el producto interno de Sobolev base:
\begin{equation}
    \ip{p}{q}_1 = \int_{S_R(C)} p(z)\overline{q(z)} d\mu(z) +  p'(a)\overline{q'(a)}
\end{equation}

Consideramos una transformación afín biyectiva del plano complejo:
$$ T(z) = \alpha z + \beta, \quad \text{con } \alpha \in \C \setminus \{0\}, \beta \in \C $$

Definimos un segundo sistema $(\mathcal{P}_2, \ip{\cdot}{\cdot}_2)$ transformado por $T$:

\begin{itemize}
    \item Nuevo Centro: $C' = \alpha C + \beta$.
    \item Nuevo Radio: $R' = |\alpha| R$.
    \item Nuevo Átomo: $a' = \alpha a + \beta$.
\end{itemize}

\begin{teorema}[Invarianza Espectral Afín]
Si $_1 = \{\xi_1, \dots, \xi_n\}$ es el conjunto de ceros del $n$-ésimo polinomio ortogonal del sistema original, entonces el conjunto de ceros del sistema transformado es:
$$ _2 = \{ \alpha \xi_1 + \beta, \dots, \alpha \xi_n + \beta \} $$
\end{teorema}

\begin{proof}
\textbf{1. Relación entre Polinomios (Demostración por Isometría):}

Sea $\{\varphi_n(z)\}_{n=0}^\infty$ la familia de polinomios ortogonales mónicos con respecto al producto interno $\ip{\cdot}{\cdot}_1$. Proponemos que los polinomios ortogonales del sistema transformado, denotados por $\{\psi_n(w)\}_{n=0}^\infty$, están dados por la composición con la transformación inversa $z = T^{-1}(w) = \frac{w-\beta}{\alpha}$. Definimos:
$$ \Psi_n(w) \coloneqq \varphi_n\left( \frac{w-\beta}{\alpha} \right) $$
Para demostrar que $\Psi_n$ es, salvo una constante multiplicativa, el $n$-ésimo polinomio ortogonal en el sistema 2, verificamos la ortogonalidad $\ip{\Psi_n}{\Psi_m}_2$ para $n \neq m$.

El producto interno en el sistema transformado es:
$$ \ip{P}{Q}_2 = \int_{S_{R'}} P(w)\overline{Q(w)} d\tilde{\mu}(w) +  P'(a')\overline{Q'(a')} $$
Sustituimos $P = \Psi_n$ y $Q = \Psi_m$. Realizamos el cambio de variable $w = \alpha z + \beta$, lo cual implica $dw = \alpha dz$. Notamos que el dominio de integración $S_{R'}$ (radio $R' = |\alpha|R$) se mapea biyectivamente a $S_R$.
\begin{itemize}
    \item \textbf{Término Integral:} Asumiendo que la medida está normalizada por la longitud de arco ($d\tilde{\mu}(w) = \frac{|dw|}{2\pi R'}$), tenemos:
    $$ \int_{S_{R'}} \Psi_n(w)\overline{\Psi_m(w)} \frac{|dw|}{2\pi |\alpha|R} = \int_{S_R} \varphi_n(z)\overline{\varphi_m(z)} \frac{|\alpha||dz|}{2\pi |\alpha|R} = \int_{S_R} \varphi_n(z)\overline{\varphi_m(z)} d\mu(z) $$
    
    \item \textbf{Término Atómico:} Usamos la regla de la cadena para la derivada. Dado $\Psi_n(w) = \varphi_n(\frac{w-\beta}{\alpha})$:
    $$ \frac{d}{dw}\Psi_n(w) = \varphi_n'\left(\frac{w-\beta}{\alpha}\right) \cdot \frac{d}{dw}\left(\frac{w-\beta}{\alpha}\right) = \varphi_n'(z) \cdot \frac{1}{\alpha} $$
    Evaluando en el átomo transformado $a' = \alpha a + \beta$ (que corresponde a $z=a$):
    $$  \Psi_n'(a')\overline{\Psi_m'(a')} =  \left( \frac{\varphi_n'(a)}{\alpha} \right) \overline{\left( \frac{\varphi_m'(a)}{\alpha} \right)} = \frac{1}{|\alpha|^2} \left(  \varphi_n'(a)\overline{\varphi_m'(a)} \right) $$
\end{itemize}

Si ajustamos el peso del átomo en el sistema 2 para compensar el escalado (o aceptamos que la ortogonalidad se mantiene salvo factores de peso constantes que no alteran la base), obtenemos:
$$ \ip{\Psi_n}{\Psi_m}_2 \propto \ip{\varphi_n}{\varphi_m}_1 = \delta_{nm} \norm{\varphi_n}_1^2 $$
Dado que $\Psi_n(w)$ es un polinomio de grado $n$ y es ortogonal a todos los polinomios de grado menor (por la isometría del producto interno), por la unicidad de los polinomios ortogonales se concluye que:
$$ \psi_n(w) = C_n \cdot \varphi_n\left( \frac{w-\beta}{\alpha} \right) $$
donde $C_n$ es una constante de normalización.

\textbf{2. Análisis del Término de la Derivada (Regla de la Cadena):}

Analizamos el comportamiento del funcional lineal asociado a la derivada bajo la transformación biyectiva $w = T(z) = \alpha z + \beta$.
Definimos los polinomios en el espacio transformado $\mathcal{H}_2$ como la composición con la transformación inversa:
$$ \Psi_n(w) = (\varphi_n \circ T^{-1})(w) = \varphi_n\left( \frac{w-\beta}{\alpha} \right) $$
Para evaluar el producto interno en $\mathcal{P}_2$, es necesario calcular la derivada compleja $\frac{d}{dw}\Psi_n(w)$. Aplicamos la regla de la cadena:
$$ \frac{d}{dw}\Psi_n(w) = \varphi_n'\left( \frac{w-\beta}{\alpha} \right) \cdot \frac{d}{dw}\left( \frac{w-\beta}{\alpha} \right) $$
Dado que la derivada del argumento lineal es constante ($\frac{d}{dw}(\frac{w}{\alpha} - \frac{\beta}{\alpha}) = \frac{1}{\alpha}$), obtenemos la relación fundamental entre las derivadas de ambos espacios:
\begin{equation}
    \Psi_n'(w) = \frac{1}{\alpha} \varphi_n'(z) \quad \text{donde } z = T^{-1}(w)
\end{equation}

Procedemos a evaluar el término atómico del producto interno $\ip{\cdot}{\cdot}_2$ en el punto transformado $a' = \alpha a + \beta$. Observamos que $T^{-1}(a') = a$.
Sustituyendo la relación derivada:
\begin{align*}
     \Psi_n'(a') \overline{\Psi_m'(a')} &=  \left[ \frac{1}{\alpha} \varphi_n'(a) \right] \overline{\left[ \frac{1}{\alpha} \varphi_m'(a) \right]} \\
    &=  \left( \frac{1}{\alpha} \right) \overline{\left( \frac{1}{\alpha} \right)} \varphi_n'(a) \overline{\varphi_m'(a)} \\
    &= \frac{1}{\alpha \overline{\alpha}} \varphi_n'(a) \overline{\varphi_m'(a)} \\
    &= \frac{1}{|\alpha|^2} \left(  \varphi_n'(a) \overline{\varphi_m'(a)} \right)
\end{align*}

\textit{Interpretación:} La transformación induce un factor de escala real positivo $k = |\alpha|^{-2}$ en el término de Sobolev.
\begin{itemize}
    \item Si la transformación es una isometría rígida (traslación o rotación pura), entonces $|\alpha|=1$, el factor es 1, y el término atómico es invariante.
    \item Si la transformación incluye una homotecia ($|\alpha| \neq 1$), el peso efectivo del átomo se escala, preservando la estructura ortogonal salvo por una constante de normalización global.
\end{itemize}

\textbf{3. Mapeo de los Ceros (Invarianza Espectral):}

Sea $\Lambda_1 = \{\xi_1, \xi_2, \dots, \xi_n\} \subset \C$ el conjunto de ceros (espectro) del $n$-ésimo polinomio ortogonal $\varphi_n$ del sistema original. Es decir, $\xi \in _1 \iff \varphi_n(\xi) = 0$.

Buscamos caracterizar el conjunto $\Lambda_2$ de ceros del polinomio transformado $\psi_n(w)$. Basándonos en la relación estructural demostrada anteriormente:
$$ \psi_n(w) = C_n \cdot \varphi_n\left( T^{-1}(w) \right) $$
donde $T^{-1}(w) = \frac{w - \beta}{\alpha}$ es la transformación inversa y $C_n \neq 0$ es una constante de normalización.

Analizamos la condición de anulación para un $\eta \in \C$:
\begin{align*}
    \eta \in \Lambda_2 &\iff \psi_n(\eta) = 0 \\
    &\iff C_n \cdot \varphi_n\left( \frac{\eta - \beta}{\alpha} \right) = 0 \\
    &\iff \varphi_n\left( \frac{\eta - \beta}{\alpha} \right) = 0 \quad (\text{pues } C_n \neq 0)
\end{align*}

Por definición de $\Lambda_1$, la función $\varphi_n$ se anula si y solo si su argumento pertenece al conjunto de sus raíces. Por lo tanto:
$$ \frac{\eta - \beta}{\alpha} \in \Lambda_1 $$
Esto implica que debe existir algún $\xi_k \in \Lambda_1$ tal que:
$$ \frac{\eta - \beta}{\alpha} = \xi_k \implies \eta = \alpha \xi_k + \beta $$

\textit{Conclusión:} El mapeo $T(z) = \alpha z + \beta$ establece una biyección entre los ceros de ambos sistemas. El espectro del sistema transformado es la imagen directa del espectro original bajo la transformación afín:
$$ \Lambda_2 = T(\Lambda_1) = \{ \alpha \xi + \beta : \xi \in \Lambda_1 \} $$
Esto demuestra que la configuración relativa de los ceros es un invariante geométrico rígido (salvo rotación y escala) respecto a la geometría del soporte de la medida.
\end{proof}

\begin{proposicion}[Transformación Afín de la Matriz de Hessenberg]
Sean $\mathbf{D}^{(1)}$ y $\mathbf{D}^{(2)}$ las matrices de representación del operador de multiplicación por la variable independiente en las bases ortonormales $\{\varphi_n\}_{n=0}^\infty$ (sistema original) y $\{\psi_n\}_{n=0}^\infty$ (sistema transformado), respectivamente. Si la transformación geométrica está dada por $w = \alpha z + \beta$, entonces:
\begin{equation}
    \mathbf{D}^{(2)} = \alpha \mathbf{D}^{(1)} + \beta \mathbf{I}
\end{equation}
\end{proposicion}

\begin{proof}
Los coeficientes de la matriz de Hessenberg en el sistema transformado se definen mediante el producto interno $\ip{\cdot}{\cdot}_2$ como:
$$ d_{ij}^{(2)} = \ip{w \psi_j}{\psi_i}_2 $$
Utilizamos el isomorfismo isométrico $\Phi: \Poly \to \Poly$ establecido anteriormente, el cual mapea la base ortonormal original a la transformada: $\Phi(\varphi_n) = \psi_n$. Debido a la invarianza del producto interno bajo este mapeo ($\ip{\Phi(p)}{\Phi(q)}_2 = \ip{p}{q}_1$), podemos llevar el cálculo al espacio original invirtiendo la transformación de la variable $w$:
$$ w \psi_j(w) = (\alpha z + \beta) \Phi(\varphi_j)(w) = \Phi( (\alpha z + \beta)\varphi_j )(w) $$
Sustituyendo esto en la definición de los coeficientes matriciales:
\begin{align*}
    d_{ij}^{(2)} &= \ip{\Phi( (\alpha z + \beta)\varphi_j )}{\Phi(\varphi_i)}_2 \\
    &= \ip{(\alpha z + \beta)\varphi_j}{\varphi_i}_1 \quad (\text{por isometría})
\end{align*}
Aplicamos la linealidad del producto interno en el primer argumento:
\begin{align*}
    d_{ij}^{(2)} &= \ip{\alpha z \varphi_j}{\varphi_i}_1 + \ip{\beta \varphi_j}{\varphi_i}_1 \\
    &= \alpha \underbrace{\ip{z \varphi_j}{\varphi_i}_1}_{d_{ij}^{(1)}} + \beta \underbrace{\ip{\varphi_j}{\varphi_i}_1}_{\delta_{ij}}
\end{align*}
Donde hemos utilizado que la base $\{\varphi_n\}$ es ortonormal en el sistema 1 ($\ip{\varphi_j}{\varphi_i}_1 = \delta_{ij}$, la delta de Kronecker).
En términos matriciales, la ecuación $d_{ij}^{(2)} = \alpha d_{ij}^{(1)} + \beta \delta_{ij}$ corresponde exactamente a:
$$ \mathbf{D}^{(2)} = \alpha \mathbf{D}^{(1)} + \beta \mathbf{I} $$
Esto confirma que la matriz del operador sufre la misma transformación afín que el plano complejo subyacente.
\end{proof}

La estructura espectral de los polinomios de Sobolev es rígida bajo transformaciones afines.
\begin{itemize}
    \item Una \textbf{traslación} ($\beta$) desplaza el espectro.
    \item Una \textbf{homotecia/rotación} ($\alpha$) escala y rota el espectro.
\end{itemize}

\end{document}